
\documentclass{article}
\usepackage{colortbl}
\usepackage{makecell}
\usepackage{multirow}
\usepackage{supertabular}

\begin{document}

\newcounter{utterance}

\centering \large Interaction Transcript for game `eqbench', experiment `eq\_bench', episode 156 with reflexion{-}Qwen3.5{-}9B.
\vspace{24pt}

{ \footnotesize  \setcounter{utterance}{1}
\setlength{\tabcolsep}{0pt}
\begin{supertabular}{c@{$\;$}|p{.15\linewidth}@{}p{.15\linewidth}p{.15\linewidth}p{.15\linewidth}p{.15\linewidth}p{.15\linewidth}}
    \# & \multicolumn{2}{c}{Player} && \multicolumn{2}{c}{Game Master} \\
    \hline

    \theutterance \stepcounter{utterance}  
    & & & \multicolumn{4}{p{0.6\linewidth}}{
        \cellcolor[rgb]{0.9,0.9,0.9}{
            \makecell[{{p{\linewidth}}}]{
                \texttt{\tiny{[P1$\langle$GM]}}
                \\ 
\texttt{Your task is to predict the likely emotional responses of a character in this dialogue:} \\
\\ 
\texttt{Eilis: Remy, you're late.} \\
\texttt{Remy: And you're as punctual as ever, Eilis. Still chasing shadows, are we?} \\
\texttt{Eilis: Better shadows than illusions of grandeur, Remy.} \\
\texttt{Remy: Ah, the sting of your wit. I've missed that. } \\
\texttt{Eilis: And I've missed your incessant need to prove something.} \\
\texttt{Remy: Is that so? Well, I suppose we're both here for the same reason then.} \\
\texttt{Eilis: Don't flatter yourself, Remy. Our reasons couldn't be more different.} \\
\texttt{{[}End dialogue{]}} \\
\\ 
\texttt{At the end of this dialogue, Remy would feel...} \\
\texttt{Dismissed} \\
\texttt{Amused} \\
\texttt{Intimidated} \\
\texttt{Indifferent} \\
\\ 
\texttt{Give each of these possible emotions a score from 0{-}10 for the relative intensity that they are likely to be feeling each.} \\
\\ 
\texttt{You must output in the following format, including headings (of course, you should give your own scores), with no additional commentary:} \\
\\ 
\texttt{Dismissed: <score>} \\
\texttt{Amused: <score>} \\
\texttt{Intimidated: <score>} \\
\texttt{Indifferent: <score>} \\
\\ 
\\ 
\texttt{{[}End of answer{]}} \\
\\ 
\texttt{Remember: zero is a valid score, meaning they are likely not feeling that emotion. You must score at least one emotion > 0.} \\
\\ 
\texttt{Your answer:} \\
            }
        }
    }
     \\ \\

    \theutterance \stepcounter{utterance}  
    & \multicolumn{4}{p{0.6\linewidth}}{
        \cellcolor[rgb]{0.561,0.737,0.561}{
            \makecell[{{p{\linewidth}}}]{
                \texttt{\tiny{[P1$\rangle$GM]}}
                \texttt{Dismissed: 8} \\
\texttt{Amused: 2} \\
\texttt{Intimidated: 3} \\
\texttt{Indifferent: 1} \\
            }
        }
    }
    & & \\ \\

    \theutterance \stepcounter{utterance}  
    & & & \multicolumn{2}{p{0.3\linewidth}}{
        \cellcolor[rgb]{0.9,0.9,0.9}{
            \makecell[{{p{\linewidth}}}]{
                \texttt{\tiny{[GM$|$GM]}}
                \texttt{Dismissed: 8} \\
\texttt{Amused: 2} \\
\texttt{Intimidated: 3} \\
\texttt{Indifferent: 1} \\
            }
        }
    }
    & & \\ \\

    \theutterance \stepcounter{utterance}  
    & & & \multicolumn{2}{p{0.3\linewidth}}{
        \cellcolor[rgb]{0.9,0.9,0.9}{
            \makecell[{{p{\linewidth}}}]{
                \texttt{\tiny{[GM$|$GM]}}
                \texttt{Dismissed: 7} \\
\texttt{Amused: 6} \\
\texttt{Intimidated: 0} \\
\texttt{Indifferent: 3} \\
            }
        }
    }
    & & \\ \\

    \theutterance \stepcounter{utterance}  
    & & & \multicolumn{2}{p{0.3\linewidth}}{
        \cellcolor[rgb]{0.9,0.9,0.9}{
            \makecell[{{p{\linewidth}}}]{
                \texttt{\tiny{[GM$|$GM]}}
                \texttt{game\_result = LOSE} \\
            }
        }
    }
    & & \\ \\

    \theutterance \stepcounter{utterance}  
    & & \multicolumn{4}{p{0.6\linewidth}}{
        \cellcolor[rgb]{0.95,0.95,0.95}{
            \makecell[{{p{\linewidth}}}]{
                \texttt{\tiny{[P1$\leftrightarrow$P]}}
                \texttt{Post{-}episode Reflexion (not sent to GM):} \\
\texttt{New plan: I failed because I assigned high scores to "Dismissed" and "Intimidated" based on the surface{-}level conflict, ignoring the underlying dynamic of mutual recognition and shared purpose. Remy's line "I've missed that" and the agreement "we're both here for the same reason" indicate he is not feeling dismissed or threatened; rather, he is engaged and finding the interaction stimulating. My strategy should have prioritized "Amused" as the primary emotion, reflecting his enjoyment of the wit and the intellectual sparring, while keeping "Dismissed" and "Intimidated" at zero. Next time, I will re{-}evaluate the dialogue for signs of shared connection and intellectual play, assigning the highest score to "Amused" (e.g., 9) and the lowest scores to negative emotions like "Dismissed" (0) and "Intimidated" (0), ensuring the scores accurately reflect the shift from conflict to camaraderie.} \\
            }
        }
    }
    & \\ \\

\end{supertabular}
}

\end{document}
